\chapter{数学符号和公式}

\section{数学符号}

中文论文的数学符号默认遵循 GB/T 3102.11—1993《物理科学和技术中使用的数学符号》
\footnote{原 GB 3102.11—1993，自 2017 年 3 月 23 日起，该标准转为推荐性标准。}。
该标准参照采纳 ISO 31-11:1992 \footnote{目前已更新为 ISO 80000-2:2019。}，
但是与 \TeX{} 默认的美国数学学会（AMS）的符号习惯有所区别。
具体地来说主要有以下差异：
\begin{enumerate}
  \item 大写希腊字母默认为斜体，如
    \begin{equation*}
      \Gamma \Delta \Theta \Lambda \Xi \Pi \Sigma \Upsilon \Phi \Psi \Omega.
    \end{equation*}
  \item 小于等于号和大于等于号使用倾斜的字形 $\le$、$\ge$。
  \item 积分号使用正体，比如 $\int$、$\oint$。
  \item 偏微分符号 $\partial$ 使用正体。
  \item 省略号按照中文的习惯固定居中，比如
    \begin{equation*}
      1, 2, \dots, n \quad 1 + 2 + \dots + n.
    \end{equation*}
  \item 实部 $\Re$ 和虚部 $\Im$ 的字体使用罗马体。
\end{enumerate}

以上数学符号样式的差异可以在模板中统一设置。
另外国标还有一些与 AMS 不同的符号使用习惯，需要用户在写作时进行处理：
\begin{enumerate}
  \item 数学常数和特殊函数名用正体，如
    \begin{equation*}
      \pi = 3.14\dots; \quad
      i^2 = -1; \quad
      e = \lim_{n \to \infty} \left( 1 + \frac{1}{n} \right)^n.
    \end{equation*}
  \item 微分号使用正体，比如 $\mathrm{d} y / \mathrm{d} x$。
  \item 向量、矩阵和张量用粗斜体，如 $\mathbf{x}$、$\mathbf{\Sigma}$。
  \item 自然对数用 $\ln x$ 不用 $\log x$。
\end{enumerate}



\section{数学公式}

数学公式可以使用 \texttt{equation} 和 \texttt{equation*} 环境。
注意数学公式的引用应前后带括号，通常使用 \texttt{\textbackslash eqref} 命令，比如式~\eqref{eq:example}。
\begin{equation}
  \frac{1}{2 \pi i} \int_\gamma f = \sum_{k=1}^m n(\gamma; a_k) \mathscr{R}(f; a_k).
  \label{eq:example}
\end{equation}

多行公式尽可能在"="处对齐，推荐使用 \texttt{align} 环境。
\begin{align}
  a & = b + c + d + e \\
    & = f + g
\end{align}



\section{数学定理}

定理环境的格式可以使用 \texttt{amsthm} 宏包配置。
模板已经配置了 \texttt{theorem}、\texttt{proof} 等环境。

\begin{theorem}[Lindeberg--Lévy 中心极限定理]
  设随机变量 $X_1, X_2, \dots, X_n$ 独立同分布， 且具有期望 $\mu$ 和有限的方差 $\sigma^2 \ne 0$，
  记 $\bar{X}_n = \frac{1}{n} \sum_{i=1}^n X_i$，则
  \begin{equation}
    \lim_{n \to \infty} P \left(\frac{\sqrt{n} \left( \bar{X}_n - \mu \right)}{\sigma} \le z \right) = \Phi(z),
  \end{equation}
  其中 $\Phi(z)$ 是标准正态分布的分布函数。
\end{theorem}
\begin{proof}
  显然成立。
\end{proof}
